% Part:sets-functions-relations
% Chapter: size-of-sets
% Section: enumerations-alt

\documentclass[../../../include/open-logic-section]{subfiles}

\begin{document}

\olfileid{sfr}{siz}{enm-alt}

\olsection{Enumerasi lan Himpunan kang \usetoken{S}{enumerable}}

\begin{editorial}
  Bagean iki netepake enumerasi minangka bijeksi karo (perangan wiwitan)
  $\Nat$, kaya kang ditindakake ing teori himpunan. Mula definisine
  rada beda karo ing \olref[enm]{sec}, lan mbaleni kabeh
  tuladha ing kana. Bagean iki uga rada luwih ringkes tinimbang
  bagean mau.
\end{editorial}

Himpunan winates bisa ditemtokake mung kanthi ngenumerasi
!!{element}sne. Iki kang ditindakake nalika netepake himpunan kaya mangkene:
\[
  A = \{a_1, a_2, \ldots, a_n\}.
\]
Yen !!{element}s $a_1$, \dots, $a_n$ kabeh beda,
iki menehi !!a{bijection} antarane $A$ lan $n$ wilangan natural
kapisan, yaiku $0$, \dots, $n-1$. Kosok baline, amarga saben
himpunan winates mung duwe !!{element}s kang cacahe winates,
saben himpunan winates bisa digawe korespondensi kaya mangkono.
Kanthi tembung liya, yen $A$ winates, ana
!!a{bijection} antarane $A$ lan $\{0, \dots, n-1\}$,
kanthi $n$ minangka cacah !!{element}s saka~$A$.

Yen kita ngakoni jinis himpunan tanpa wates tartamtu, kita uga
ngakoni yen sawetara himpunan tanpa wates bisa dienumerasi.
Iki bisa ditegesi kanthi pas: himpunan tanpa wates dienumerasi
dening !!a{bijection} antarane himpunan iku lan kabeh~$\Nat$.

\begin{defn}[Enumerasi, miturut teori himpunan]
\emph{Enumerasi} himpunan $A$ iku !!a{bijection} kang jangkauane
$A$ lan domaine himpunan wiwitan wilangan natural $\{0,
1, \ldots, n\}$ {utawa} kabeh himpunan wilangan natural~$\Nat$.
\end{defn}

\begin{explain}
Panganggone tembung \emph{enumerasi} iki duwe dhasar intuisi.
Ngandharake yen kita wis ngenumerasi himpunan $A$ ateges
ana !!a{bijection} $f$ kang ngidini kita ngetung anggota
himpunan $A$ siji-siji. Anggota kaping $0$ yaiku $f(0)$,
anggota kapisan $f(1)$, \ldots anggota kaping $n$ yaiku $f(n)$\ldots.
\footnote{Pancen, kita ngetung wiwit $0$. Mesthi wae kita uga
bisa miwiti saka~$1$. Ora bakal ana bedane kang wigati.
Kita mung kudu ngganti~$\Nat$ nganggo~$\PosInt$.}
Dhasar gagasan iki bisa luwih cetha kanthi tambahan:
\end{explain}

\begin{defn}
  \ollabel{defn:enumerable}
  Himpunan~$A$ iku !!{enumerable} yen lan mung yen $A = \emptyset$
  utawa ana enumerasi saka~$A$. Kita nyebut $A$ iku !!{nonenumerable}
  yen lan mung yen $A$ ora !!{enumerable}.
\end{defn}

\begin{explain}
Dadi himpunan iku !!{enumerable} yen lan mung yen kosong utawa kita
bisa nggunakake enumerasi kanggo ngetung !!{element}sne siji-siji.
\end{explain}

\begin{ex}
Fungsi kang ngenumerasi wilangan natural yaiku fungsi identitas
$\Id{\Nat} \colon \Nat \to \Nat$ kang diwenehake dening $\Id{\Nat}(n) = n$.
Fungsi kang ngenumerasi wilangan natural \emph{positif}, yaiku $\Nat^+ =
\Nat \setminus \{0\}$, yaiku fungsi $g(n) = n + 1$,
utawa fungsi panerus.
\end{ex}

\begin{prob}
Tuduhna yen himpunan $A$ iku !!{enumerable} yen lan mung yen $A = \emptyset$
utawa ana !!a{surjection} $f\colon \Nat \to A$. Tuduhna yen $A$
iku !!{enumerable} yen lan mung yen ana !!a{injection} $g\colon A \to \Nat$.
\end{prob}

\begin{ex}
Fungsi $f\colon \Nat \to \Nat$ lan $g \colon \Nat \to \Nat$
kang diwenehake dening
\begin{align*}
  f(n) & = 2n \text{ lan}\\
  g(n) & = 2n+1
\end{align*}
ngenumerasi wilangan natural genep lan wilangan natural ganjil.
Nanging loro-lorone ora !!{surjective}, mula loro-lorone dudu
enumerasi saka $\Nat$.
\end{ex}

\begin{prob}
Tetepna enumerasi wilangan kuadrat $1$, $4$, $9$, $16$, \dots
\end{prob}

\begin{ex}
Upamakna $\lceil x \rceil$ iku \emph{fungsi pambunderan munggah},
kang mbunderake $x$ munggah menyang wilangan wutuh kang paling cedhak.
Mula fungsi $f \colon \Nat \to \Int$ kang diwenehake dening:
\[
  f(n) = (-1)^{n} \left\lceil\tfrac{n}{2}\right\rceil
\]
ngenumerasi himpunan wilangan
wutuh~$\Int$ kaya mangkene:
\[
\begin{array}{c c c c c c c c}
f(0) & f(1) & f(2) & f(3) & f(4) & f(5) & f(6) & \dots \\ \\
\big\lceil \tfrac{0}{2} \big\rceil & -\big\lceil \tfrac{1}{2}\big\rceil &  \big\lceil \tfrac{2}{2} \big\rceil & -\big\lceil \tfrac{3}{2} \big\rceil & \big\lceil \tfrac{4}{2} \big\rceil  & -\big\lceil \tfrac{5}{2}\big\rceil & \big\lceil \tfrac{6}{2} \big\rceil & \dots \\ \\
0 & -1 & 1 & -2 & 2 & -3 & 3& \dots
\end{array}
\]
Gatekna carane $f$ ngasilake nilai ing $\Int$ kanthi "mlumpat"
genti-genten antarane wilangan wutuh positif lan negatif.
Kita uga bisa nganggep $f$ ditetepake kanthi kasus kaya mangkene:
\[
f(n) = \begin{cases}
  \frac{n}{2} & \text{yen $n$ genep}\\
  -\frac{n+1}{2} & \text{yen $n$ ganjil}
  \end{cases}
\]
\end{ex}

\begin{prob}
Tuduhna yen $A$ lan $B$ iku !!{enumerable}, mula $A \cup B$ uga mangkono.
\end{prob}

\begin{prob}
Tuduhna kanthi induksi tumrap $n$ yen $A_1$, $A_2$, \dots, $A_n$ kabeh
!!{enumerable}, mula $A_1 \cup \dots \cup A_n$ uga mangkono.
\end{prob}

\end{document}
